\section{Discussion and Limitations}
\textbf{Practice.} (1) Audit every attribute you will report on, and crossings where cell sizes allow; a sex-only audit passed a model with a 12$\times$ larger intersectional gap, while an audit of both marginals would not have. (2) Report DP, EO, sensitivity, PPV and calibration together, at more than one operating point. (3) Compare mitigations at equal AUROC and across penalty strengths: the EO gap is U-shaped in $\lambda$, and a threshold baseline that looks clearly better against a strong penalty only ties a weak one. (4) Do not credit a federated method for an effect a centralized control reproduces, and test it under client heterogeneity; local-count re-weighting is the one mitigation here that federation breaks. (5) Train to convergence, with a schedule matched to penalty stiffness; our own earlier under-trained runs gave different answers.

\textbf{Privacy cost of thresholds.} Fitting group thresholds in FL needs, per group, score quantiles or histograms from each client; with 150 respondents per client the younger-women cell holds about 40 respondents and about one positive, so the shared statistics can approach patient-level information, and a differentially private histogram would add noise to exactly the groups the thresholds protect. We fitted thresholds centrally and make no claim about a private federated procedure.

\textbf{Limitations.} (1) One self-reported US survey with near-IID states; federated claims under skew rest on synthetic Dirichlet partitions; two binary attributes only, so the value of crossing three or more (for example race or income, absent from the cleaned extract) is untested and is future work, and marginal sufficiency may not carry over. (2) Logistic regression only in the main text (an MLP is exploratory, \sref{sec:mlp}); FairFed, FedIDA and exponentiated gradient are not benchmarked. (3) The training schedule is checked against half the rate and twice the rounds on a subset of settings (\sref{app:conv}), not for every run. (4) Seeds are random splits of one dataset (5--20 per suite); with 10 or fewer seeds, ``not significant'' does not mean ``absent''. (5) DP is a parity audit, not a clinical objective; gaps depend on the operating point and, for offsets, on the $\lambda$ grid; $\lambda$ is swept, not tuned. (6) Thresholds use the sensitive attribute at decision time, and the cost of estimating them privately is unmeasured.

\section{Conclusion}
A sex-only audit passes models whose age$\times$sex gaps remain large, but penalizing both marginals closes them about as well as penalizing the cross. In natural data penalties cost about 0.12 AUROC and 30\% of sensitivity at $\lambda{=}8$ and 0.05 AUROC at the EO-optimal $\lambda{=}0.5$; group thresholds keep AUROC, beat a strong penalty on EO and tie a weak one, but need the sensitive attribute and are unreliable when fitted on small validation splits. Local oversampling's DP/EO shift at equal AUROC is partly (49--90\% of the EO shift) a class-rebalancing effect, is reproduced by a centralized control and is significant only in natural-size clients. Group re-weighting with local counts fails under skewed clients and global counts repair it. The practical contribution is an evaluation protocol: group-wise and marginal audits, DP and EO with sensitivity and PPV, equal accuracy over several penalty strengths and operating points, a thresholding baseline, a centralized control and convergence checks.

